\subsection{From proved capture guards to timing exceptions}
\label{sec:chain}

A PE result is a combinational function of the PE's own register, its
neighbours' registers, its configuration word and the edge registers. The
controller captures it only every $D$ clocks, so the design applies a
$D$-cycle setup and $(D-1)$-cycle hold multicycle exception to these
paths~\cite{ug903,bhasker-timing}. Such an exception is correct only if the
sources are stable for $D$ cycles before every capture; a positive slack
under the exception does not establish that premise. We document the
argument as a chain of four checkable links.

\emph{(1) Functional premise, proved on the RTL.} For the controller,
properties proved by $k$-induction on every enumerated instance
(Table~\ref{tab:formal}) state that captures are at least $D$ clocks apart,
that the first capture after a datapath reset is at least $D$ clocks later,
that configuration and edge registers change only in their receive states
or on reset, and that at least $D-1$ clocks separate such a change from the
next non-reset capture. The PE and array proofs establish that PE registers
change only on step or reset. For the block extension, analogous guards on
the operand and ownership registers are proved for $D=1$--4. The
reachable-state lemmas these inductions need are asserted and proved, never
assumed: the controller tasks have no assumptions at all, and the
integrated block controller assumes only an initial reset; it proves the
range bound ($\mathit{count}\le1024$) that the standalone sequencer and
memory tasks assume.

\emph{(2) Constraint scope.} The constraint script names every expected
source register bit---for $4\times4$, 256 PE result bits, 416 used
configuration bits, 64 west and 64 north bits, i.e.\ 800 launch registers
(833 with the block operand and ownership registers)---and the 256 PE data
inputs as endpoints. Exceptions are applied from clock pins to data pins
only, so clock-enable and reset pins keep single-cycle checks, as do RAM to
operand-register paths. A run fails if any expected register is missing
from the synthesized netlist or any unexpected one matches.

\emph{(3) Static timing.} The routed netlist must meet setup and hold with
these exceptions, and must pass checks on endpoint coverage, unconstrained
paths, synchronisers and clock-domain crossings.

\emph{(4) Residual gap.} The link between (1) and (2) is by register name:
we did not check formally that synthesis preserved the RTL register
boundaries, and the proofs cover enumerated geometries and $D$ values, not
all parameters. Metastability of the asynchronous inputs is outside both
proofs and STA~\cite{ginosar-metastability}.

Table~\ref{tab:vivado} lists the implementation results. On the current
revision the legacy $4\times4$ design closes at 10\,ns with the same slack
and resources as an archived run on an earlier snapshot, which differed only
in a top-level wrapper that is pure wiring when the extension is disabled.
A heuristic period search on that earlier snapshot found a passing
candidate at 5.859\,ns ($1000/T=170.7$\,MHz); implementing the current
revision at that period again passes every gate with identical slack and
resources. This is a constraint result for one placement, not a maximum
frequency and not a change of the board's 100\,MHz oscillator. With $D=2$, the corresponding step rate is
$f_{\mathrm{clk}}/D$, and kernel throughput is lower still by $\eta$ and,
dominantly, by the link (RQ1). The block extension closed at 10\,ns only
after two retained failing designs ($-3.580$ and $-2.867$\,ns), which led
to explicit operand and ownership registers with proved guards. The
geometry sweep uses one placement per shape, so differences in slack
between shapes are not statistically interpretable.

\begin{table*}[t]
  \centering
  \caption{Formal results on the production VHDL, re-executed for this paper
  on the current sources (SymbiYosys with Yices~2.7.0, GHDL~6.0.0 plugin,
  Yosys~0.67). PROVED: bounded base case and $k$-induction at the stated
  depth. Instances are enumerated; no result is a theorem over all
  parameters. Assertion counts are cells in the solver model.}
  \label{tab:formal}
  \footnotesize
  \renewcommand{\arraystretch}{1.12}
  \begin{tabularx}{\textwidth}{@{}>{\raggedright\arraybackslash}p{2.6cm}
      >{\raggedright\arraybackslash}X >{\raggedright\arraybackslash}p{3.6cm}
      >{\raggedright\arraybackslash}p{1.5cm} r >{\raggedright\arraybackslash}p{2.2cm}@{}}
    \headtop
    \headrow
    \thead{Family} & \thead{Instances and conclusion} & \thead{Assumptions}
      & \thead{Method} & \thead{Assert.} & \thead{Result} \\
    \headbottom
    PE datapath & All 16 opcodes, both 8-way muxes, all 16-bit operands; reset priority and hold
      & Initial reset & $k$-ind., $k=4,8$ & 1, 4 & PROVED \\
    Array wiring & $1\times1$ monolithic; $1\times2$, $2\times1$, $2\times3$, $3\times2$, $4\times4$, $1\times16$, $16\times1$ by composition
      & Initial reset; PE transitions as proved component guarantees & $k$-ind., $k=4$ & 1--16 & PROVED \\
    Legacy controller & $1\times1$ with $D=1$--4; $2\times3$, $3\times2$, $4\times4$, $1\times16$, $16\times1$ with $D=2$; capture count/spacing, guards, rejected \cmd{EXEC}, \cmd{EXEC} reply, watchdog, TX handshake, ranking functions
      & None & $k$-ind., $k=16$ & 77--78 & PROVED \\
    UART & TX and RX digital models at divisors 3 and 868
      & None & $k$-ind., $k=16$ & 21, 31 & PROVED \\
    Block sequencer & $D=1$--4: count, cadence, operand/ownership guards, result capture
      & $\mathit{count}\le1024$ (port subtype) & $k$-ind., $k=16$ & 24 & PROVED \\
    Block memories & Two $1024\times16$ banks, arbitrary address, $D=1$--4: retention and read order
      & Same & $k$-ind., $k=16$ & 31 & PROVED \\
    Frame receiver & Timeout 3 and $10^{8}$: escaping, length, CRC accept/reject, watchdog
      & None & $k$-ind., $k=8$ & 22 & PROVED \\
    Block controller & Receiver and sequencer connected, $D=1$--4 and $T=10^{8}$: effects only from validated frames, cache, exact replay without duplicated effects, sequence exhaustion
      & Initial reset & $k$-ind., $k=6$ & 86 & PROVED \\
    \midrule
    Negative controls & One false assertion and nine mutated copies (adder, step, \cmd{EXEC} byte order, north link, UART TX/RX, CRC bypass, RAM bank, early capture)
      & As above & BMC, 4--40 & -- & 10/10 FAIL as required \\
    Reachability & 18 explicit cover tasks, including a complete block transcript driver
      & Witness-only stimuli & BMC, 40--200 & -- & 18 BOUNDED\_PASS \\
    \bottomrule
  \end{tabularx}
\end{table*}


\begin{table}[t]
  \centering
  \caption{Routed implementation results, Vivado 2026.1,
  \texttt{xc7a100tcsg324-1}, $D=2$, default strategy, no floorplan, one
  placement each. All runs passed every gate (setup, hold, endpoint and
  constraint coverage, synchronisers, CDC, DRC). LUT: post-synthesis /
  post-route. Rev.: C~=~current revision; E~=~earlier snapshot whose
  top-level wrapper differs only by wiring; ``this paper'' marks runs
  executed for this revision.}
  \label{tab:vivado}
  \footnotesize
  \setlength{\tabcolsep}{3pt}
  \begin{tabular}{@{}llrrrrc@{}}
    \headtop
    \headrow
    \thead{Design} & \thead{$T$/ns} & \thead{WNS} & \thead{WHS} & \thead{LUT} & \thead{FF} & \thead{Rev.} \\
    \headbottom
    $4\times4$, this paper     & 10    & $+1.862$ & $+0.098$ & 4\,999/4\,971 & 1\,034 & C \\
    $4\times4$, this paper     & 5.859 & \VivadoFastWNS & \VivadoFastWHS & \VivadoFastLUT & \VivadoFastFF & C \\
    $4\times4$, period search  & 5.859 & $+0.177$ & $+0.122$ & --/4\,975 & 1\,034 & E \\
    $4\times4$ with blocks     & 10    & $+0.427$ & $+0.060$ & 8\,262/8\,164 & 4\,346 & C \\
    \midrule
    $1\times16$ & 10 & $+2.433$ & $+0.081$ & 5\,017/4\,999 & 1\,179 & E \\
    $2\times8$  & 10 & $+0.775$ & $+0.105$ & 5\,003/4\,974 & 1\,066 & E \\
    $4\times4$  & 10 & $+1.862$ & $+0.098$ & 4\,999/4\,971 & 1\,034 & E \\
    $8\times2$  & 10 & $+2.350$ & $+0.130$ & 5\,002/4\,975 & 1\,066 & E \\
    $16\times1$ & 10 & $+1.247$ & $+0.098$ & 5\,010/4\,993 & 1\,179 & E \\
    \bottomrule
  \end{tabular}
\end{table}

